% !TeX program = xelatex
% UTF-8; compile from this directory with XeLaTeX.
% To input in an existing XePersian document, define \FetalGrowthInNotebook.
% Source: Google Drive Notebook/fetal growth, handwritten pages 1--6.
% Original JPEGs, in page order:
% pic01.jpg = 20261013_025113440.jpg
% pic02.jpg = 20261013_025125569.jpg
% pic03.jpg = 20261013_025131497.jpg
% pic04.jpg = 20261013_025141548.jpg
% pic05.jpg = 20261013_025147113.jpg
% pic06.jpg = 20261013_025151993.jpg
\ifdefined\FetalGrowthInNotebook
  \let\FGfinish\relax
\else
  \def\FGfinish{\end{document}}
  \documentclass[12pt,a4paper]{article}
  \usepackage[margin=22mm]{geometry}
  \usepackage{amsmath,amssymb}
  \usepackage{graphicx}
  \usepackage{xepersian}
  \IfFontExistsTF{Amiri}{\settextfont{Amiri}}{\settextfont{DejaVu Sans}}
  \setlatintextfont{DejaVu Serif}
  \setlength{\parindent}{0pt}
  \setlength{\parskip}{6pt}
  \setlength{\emergencystretch}{3em}
  \begin{document}
\fi
\providecommand{\FGpage}[1]{\clearpage\section*{صفحهٔ #1}}
\providecommand{\FGmath}[1]{\lr{\(#1\)}}
\providecommand{\FGunclear}[1]{\textbf{[خوانش نامطمئن: #1]}}
\providecommand{\FGnote}[1]{\textnormal{[یادداشت رونویسی: #1]}}
\providecommand{\FGimage}[1]{%
  \IfFileExists{images/#1.jpg}{\includegraphics[width=\textwidth,height=.87\textheight,keepaspectratio]{images/#1.jpg}}%
  {\includegraphics[width=\textwidth,height=.87\textheight,keepaspectratio]{Ob&Gyn/Obstetrics/fetal growth/images/#1.jpg}}}
\providecommand{\FGheader}{\textbf{سربرگ چاپی:} \lr{Subject: Year. Month Date. ( )}؛ فیلدهای تاریخ خالی. نشان چاپی پایین صفحه: \lr{PAPCO}.}

\begin{center}\Large\bfseries\lr{Fetal Growth and Surveillance}\end{center}
رونویسی متن روی شش صفحه، به ترتیب شمارهٔ دست‌نویس. شاخه‌ها و فلش‌ها به صورت فهرست و مراحل متوالی آمده‌اند. عبارت‌ها و اعداد منبع، از جمله اختصارات و موارد خط‌خورده، حفظ شده‌اند؛ خوانش‌های نامطمئن مشخص شده‌اند. تصاویر کامل منبع در پیوست همین فایل آمده‌اند. اثر کم‌رنگ و معکوس نوشته‌های پشت کاغذ متن روی صفحه محسوب نشده است.

\FGpage{۱}
\FGheader\par
\textbf{موضوع دست‌نویس سربرگ:} \lr{Fetal Growth and Surveillance}. شمارهٔ پایین صفحه: ۱.

\begin{latin}
\begin{itemize}
\item Fetal Growth Restriction (FGR)
\item Small for Gestational Age (SGA)
\end{itemize}
\end{latin}
پروگنوز \lr{SGA} بهتر از \lr{FGR} است.

\subsection*{تخمین زمان زایمان و سن بارداری}
تخمین زمان زایمان ۴۰ هفته پس از روز شروع \lr{LMP} است.
\begin{itemize}
\item اگر \lr{cycle}های فرد منظم و ۲۸روزه باشند؛ اگر مثلاً ۳۵روزه باشند، باید ۷ روز به \lr{LMP} اضافه نمود.
\item بهترین روش تعیین سن بارداری سونوگرافی پیش از ۲۰هفتگی است.
\item دقیق‌ترین حالت: تعیین \lr{Crown Rump length} بین ۱۰–۱۴هفتگی است.
\item به علاوه \lr{HC} بین ۱۵–۲۰هفتگی.
\end{itemize}

\subsection*{\lr{SGA} و \lr{FGR}}
\begin{itemize}
\item \lr{SAG}: وزن تخمینی جنین و یا وزن تولد جنین پایین‌تر از صدک ۱۰ باشد.
  \FGnote{در این سطر اختصار به صورت \lr{SAG} نوشته شده است.}
\item \lr{FGR}: جنین به هر علتی نتوانسته به پتانسیل رشد ژنتیکی خود برسد.
\item رشد جنین در سونوگرافی‌های متوالی کمتر از حد انتظار است.
\item این حد به طور دقیق مشخص نیست، بلکه بنا به تجربهٔ فرد است.
\item نوزاد با \lr{FGR}، پایین است \lr{(ponderal Index)}، \FGmath{(BW/height)}؛ و چربی زیرپوستی کمی دارد.
\end{itemize}

\textbf{نکته:} بسیاری از \lr{SGA}ها \lr{FGR} نیستند، ولی اغلب آن‌ها صرفاً جنین‌هایی هستند که به طور ژنتیکی کوچک‌اند. برخی از جنین‌های \lr{FGR} دچار محدودیت شده‌اند ولی هنوز به زیر صدک نرسیده‌اند.

\subsection*{فاکتورهای جنینی که رشد جنین را تحت تأثیر قرار می‌دهند}
\begin{itemize}
\item ژنتیک جنین مهم‌ترین تعیین‌کنندهٔ رشد آن است.
\item این ژنتیک بیشتر از مادر به ارث می‌رسد تا پدر.
\item برخی اختلالات ساختاری در جنین.
\item برخی بیماری‌ها در مادر: \lr{CKD}.
\item برخی عفونت‌ها: توکسوپلاسموز، \lr{CMV}، روبلا، مالاریا.
\end{itemize}

\subsection*{فاکتورهای مادری که رشد جنین را تحت تأثیر قرار می‌دهند}
\begin{itemize}
\item تغییرات جزئی در رژیم غذایی تأثیر مختصری دارند، ولی گرسنگی شدید طولانی مؤثر است.
\item \FGunclear{عبارت مقایسهٔ نژادها؛ خوانش نزدیک به متن: «نژادهای معین به طور نرمال جنین‌های بزرگ‌تری از حد نرمال در افریقایی‌ها رشد می‌دهند».}
\item رسیدن اکسیژن کافی به جنین بسیار مهم است.
\item مانند ارتفاع بالا در مادر یا هایپوکسی مزمن (بیماری مادرزادی قلبی).
\end{itemize}

\FGpage{۲}
\FGheader\par
شمارهٔ پایین صفحه: ۲.

\subsection*{ادامهٔ عوامل مادری}
\begin{itemize}
\item جنین می‌تواند نسبت به هایپوکسی یک \lr{partial compensation} ایجاد کند.
\item از طریق هایپرتروفی \lr{placenta}.
\item مواد مخدر از جمله تنباکو، هروئین، کوکائین و الکل رشد جنین را مختل می‌کنند.
\item سیگار به طور متوسط \lr{BW} نوزاد را \lr{150 gr} کاهش می‌دهد.
\end{itemize}

\subsection*{فاکتورهای مرتبط با جفت که رشد جنین را تحت تأثیر قرار می‌دهند}
\begin{itemize}
\item کارکرد کافی پلاسنتا وابسته به تهاجم مناسب تروفوبلاستیک است.
\item این تهاجم در سه‌ماههٔ اول رخ می‌دهد.
\item در سه‌ماههٔ دوم یک تهاجم ثانویه نیز وجود دارد.
\item در پره‌اکلامپسی این تهاجم ثانویه مختل شده است.
\item این اختلال در تهاجم با سونو داپلر قابل تشخیص است.
\item اختلال در جریان خون با نوزاد \lr{SGA} مرتبط است و نیازمند انجام سونوگرافی سریال می‌باشد.
\item نارسایی پیشروندهٔ پلاسنتال.
\item جنین هایپوکسیک، هایپوگلیسمیک، و در برخی مواقع اسیدوتیک.
\item به طور جبرانی \lr{Hb} بالا می‌رود؛ جریان خون محیطی، روده و کبد کم شده و به سمت مغز، قلب و غدهٔ آدرنال می‌رود.
\item جنین قد و رشد مغزی نرمال دارد ولی لاغر است.
\end{itemize}

\subsection*{غربالگری و تشخیص}
\begin{itemize}
\item برخی از ویژگی‌های ریسک فاکتور \lr{SGA} در اولین ویزیت مادر باردار قابل شناسایی هستند، ولی برخی در طول بارداری نمایان می‌شوند.
  \FGnote{یک \lr{SGA} پیش از «ریسک فاکتور» خط خورده است.}
\end{itemize}

\textbf{ریسک فاکتورهای \lr{SGA}:}
\begin{itemize}
\item شرح حال مادر: سن بالاتر از ۴۰ سال، سیگاری بیش از ۱۰ سال، مصرف کوکائین، سابقهٔ قبلی \lr{SGA} در نوزاد قبلی، سابقهٔ قبلی مرده‌زایی، \lr{HTN} مزمن، \lr{DM} و بیماری عروق، اختلال کلیوی، \lr{APS}.
  \FGnote{عبارت سیگار در منبع به صورت «سیگاری \FGmath{10<} سال» آمده است.}
\item بارداری فعلی: تهدید به سقط، پره‌اکلامپسی، پلاسنتال ابراپشن، خونریزی غیرقابل توجیهٔ انتپارتوم، نتایج غیرطبیعی شریان رحمی در داپلر، سطح پایین \lr{PAPP-A}، رودهٔ جنین در سونو هایپراکو باشد.
  \FGunclear{عبارت داخل پرانتز پس از انتپارتوم: «جفت‌ها کنار آمده باشند»؛ احتمال خوانش «جفت‌ها کنار آمده باشه».}
\item وجود ریسک فاکتورهای بالا.
\item سونوگرافی سریال از هفتهٔ ۲۶ تا ۲۸ برای تشخیص \lr{fetal growth}.
\end{itemize}
\textbf{نوشتهٔ زیر خط پایین صفحه:}
\begin{latin}
Plasma Associated Placenta Protein-A
\end{latin}

\FGpage{۳}
\FGheader\par
شمارهٔ پایین صفحه: ۳.

\subsection*{معاینهٔ بالینی}
\begin{itemize}
\item تخمین وزن جنین با معاینه غیرقابل اطمینان است.
\item فوندوس طی هفته‌های ۲۰–۲۴ به حدود ناف می‌رسد.
\item فوندوس در حدود هفتهٔ ۳۶ به حدود زایفوئید می‌رسد.
\item از هفتهٔ ۲۶ در هر ویزیت روتین باید ارتفاع فوندوس اندازه‌گیری شود.
\item از سمفیز پوبیس: \lr{SFH (Symphysis Fundal Height)}.
\item تقریباً معادل سن بارداری است.
\item \lr{SFH} زیر صدک ۱۰ هم، یا اگر در اندازه‌گیری سریال افزایش نداشته باشد، اندیکاسیون انجام سونوگرافی است.
\end{itemize}

\subsection*{سونوگرافی}
\begin{itemize}
\item برای تأیید و یا رد تشخیص \lr{SGA} یا \lr{FGR} استفاده می‌شود.
\item تشخیص قطعی تنها پس از تولد ممکن است.
\item مواردی که باید گزارش شوند: \lr{HC, HC, FL} (طول فمور)، \lr{EFW}.
  \FGnote{خوانش اختصار دوم میان \lr{HC} و \lr{AC} نامطمئن است.}
\item \lr{AC} برای تشخیص \lr{SGA} از بقیه حساس‌تر است.
\item زیر صدک ۱۰، حساسیت ۸۰ درصد دارد.
\item خطای \lr{EFW}، \FGmath{\pm15\%} است.
\end{itemize}

\subsection*{\lr{SGA} یا \lr{FGR}؟}
\begin{itemize}
\item سؤال مهمی است اما همیشه نمی‌توانیم پاسخ قطعی آن بدهیم.
\item جنین‌های زیر صدک ۱۰، ۵۰–۷۰ درصد \lr{SGA} هستند؛ \lr{FGR} نیستند.
\item افزایش ریسک مرده‌زایی، هایپوکسی و عوارض مربوط به \lr{FGR} است.
\item تشخیص \lr{SGA} سرراست است ولی در \lr{FGR} ساده نیست.
\end{itemize}

\textbf{کرایتریا تشخیص قطعی \lr{FGR}:} دو مورد زیر در منبع با آکولاد به هم متصل شده‌اند:
\begin{itemize}
\item افت حداقل ۵۰ صدک در \lr{EFW}، و
\item افزایش مقاومت شریان آمبیلیکال در داپلر.
\end{itemize}
یا یکی از معیارهای بیومتریک زیر صدک ۳ باشد.

زمانی که \lr{SGA} یا \lr{FGR} تشخیص داده می‌شود، باید سایر پارامترهای جنین نیز ارزیابی شوند.

\subsection*{کاربرد بالینی}
جنین با \lr{AC} یا \lr{EFW} زیر صدک ۱۰ باید \lr{Close Observation} شود؛ خواه این جنین \lr{SGA} یا \lr{FGR} یا هر دو باشد.

\FGpage{۴}
\FGheader\par
شمارهٔ پایین صفحه: ۴.

\subsection*{مدیریت بالینی}
\begin{itemize}
\item اصل اصلی: مانیتور کردن جنین و زایمان در زمان مناسب.
\item هیچ درمان خاصی اثبات‌شده نیست.
\end{itemize}

\textbf{گزینه‌ها برای مانیتور کردن جنین:}
\begin{latin}
\begin{itemize}
\item FMCC: Fetal Movement Counting Chart
\item FCTG: Fetal Cardiotocography
\item Doppler blood flow studies \(\longrightarrow\) Most Valuable
\item Amniotic Fluid Volume
\end{itemize}
\end{latin}

\subsection*{کاربرد بالینی}
\begin{enumerate}
\item مانیتورینگ سریال با سونوگرافی و/یا \lr{SFH}.
\item تأیید آنکه جنین \lr{SGA} یا \lr{FGR} است.
\item رد ابنورمالیتی‌های جنین (مگر آنکه پیش از این در هفتهٔ ۲۰ هم انجام شده باشد).
\item تشخیص اینکه جنین \lr{SGA} است یا \lr{FGR}: انجام سونوگرافی داپلر.
\item مانیتورینگ با داپلر و \lr{CTG}.
\item اگر \FGmath{GA<34} است و قصد زایمان دارید، کورتیکواستروئید بدهید.
\item اگر احتمال \lr{fetal demise} وجود دارد، زایمان را انجام بدهید.
\end{enumerate}

\subsection*{مانیتورینگ حرکات جنین}
\begin{itemize}
\item یک جنین کمتر تغذیه‌شده برای حفظ انرژی حرکات خود را کم می‌کند.
\item هر کاهش ناگهانی در تعداد حرکات جنین و الگوی آن ممکن است مهم باشد.
\end{itemize}

\subsection*{کاردیوتوکوگرافی \lr{(CTG)}}
\begin{itemize}
\item وضعیت سلامت جنین را در لحظه مشخص می‌کند ولی در \lr{long-term} ارزش کمتری دارد.
\item \lr{CTG} طی بارداری با \lr{CTG} حین زایمان تفسیر متفاوتی دارد.
\item حین بارداری، وجود \lr{Acceleration} و نبود \lr{Deceleration} مورد انتظار است.
\item طی زایمان وجود \lr{Deceleration} می‌تواند نرمال باشد.
\item کرایتریاهای بررسی \lr{CTG}: \lr{The Dawes Redman Criteria}.
\item استفادهٔ روتین از \lr{CTG} با بهبود پیامد بالینی همراهی ندارد.
\item به همین دلیل تنها زمانی استفاده می‌شود که اندیکاسیون وجود داشته باشد، مثل \lr{RFM}.
  \FGnote{کلمهٔ «کرایتریا» پیش از «اندیکاسیون» خط خورده است.}
\end{itemize}

\FGpage{۵}
\FGheader\par
شمارهٔ پایین صفحه: ۵.

\subsection*{ارزیابی حجم مایع آمنیوتیک}
\begin{itemize}
\item بازتوزیع حجم مایعات بدن در جنین طی \lr{FGR} منجر به کاهش جریان خون کلیهٔ جنین و به دنبال آن کاهش ادرار جنین و کاهش حجم مایع آمنیوتیک می‌شود.
\item وجود الیگوهیدرآمنیوس در نبود \lr{ROM} یافتهٔ کاراکتریستیک \lr{FGR} در سونوگرافی است.
\item توصیف مایع آمنیوتیک در سونوگرافی:
  \begin{itemize}
  \item \lr{MVP: Max Vertical Pool}.
  \item \lr{AFI: Amniotic Fluid Index}؛ مجموع حجم مایع در ۴ کوادران.
  \end{itemize}
\end{itemize}

\subsection*{سونوگرافی داپلر}
\begin{itemize}
\item یک \lr{waveform} طبیعی بیانگر این است که جنین \lr{FGR} نیست.
\item کاهش یا از دست رفتن \lr{End-diastolic flow} بیانگر \lr{high risk} بودن جنین برای هایپوکسی است.
\item فقدان جریان انتهای دیاستولی \lr{(AEDF)} مطرح‌کنندهٔ خطر مرگ بالا در جنین \lr{FGR} است.
\item بررسی جریان خون مغزی جنین نیز کمک‌کننده است، خصوصاً در جنین‌های ترم.
\item در جنین \lr{FGR}، کاهش مقاومت عروق در شریان \lr{MCA}، افزایش جریان خون مغز دیده می‌شود.
\item به دنبال تشدید هایپوکسی، به دنبال تشدید \lr{FGR}، مقاومت عروق مغزی که به طور جبرانی کاهش پیدا کرده بود مجدداً افزایش می‌یابد.
\item احتمالاً ثانویه به ادم مغزی.
\item بررسی جریان خون وریدی برای تعیین زمان زایمان کمک‌کننده است.
\item داپلر \lr{Ductus Venosus} (وریدی در کبد جنین) به طور غیرمستقیم عملکرد قلب راست جنین را ارزیابی می‌کند.
\item اگر شکل موج \lr{Ductus Venosus} غیرطبیعی باشد، قویاً مطرح‌کنندهٔ \lr{Fetal Demise} در آینده است.
\end{itemize}

\subsection*{سریال تغییرات به دنبال هایپوکسی جنینی}
ترتیب فلش‌ها در منبع:
\begin{enumerate}
\item اختلال رشد جنین.
\item شکل موج غیرطبیعی شریان آمبیلیکال.
\item افزایش جریان خون مغزی.
\item جریان غیرطبیعی \lr{DV}.
\item ریت قلبی غیرطبیعی جنین \lr{(FHR)}.
\item \lr{Fetal Demise}.
\end{enumerate}

\textbf{مارکرهایی که از چند هفته قبل از مرگ جنین رخ می‌دهند:}
\FGunclear{ابتدای جملهٔ عنوان، «مارکرهایی که از چند هفته قبل از مرگ جنین رخ می‌دهند».}
\begin{itemize}
\item افزایش \lr{UA PI} (اندکس پالساتیلیتی شریان آمبیلیکال): زودتر از همه رخ می‌دهد.
\item \lr{AEDF}.
\item الیگوهیدرآمنیوس.
\item هایپوکسی.
\end{itemize}

\FGpage{۶}
\FGheader\par
شمارهٔ پایین صفحه: ۶.

\subsection*{استراتژی کلی}
\begin{itemize}
\item سونوگرافی داپلر معیار غربالگری اولیهٔ ما می‌باشد.
\item در جنین با \lr{AEDF} یا فلو کاهش‌یافته: مانیتورینگ با داپلر و \lr{CTG}.
\item زایمان اگر \lr{CTG} غیرطبیعی و پاتولوژیک شد اندیکاسیون دارد.
\item اگر \lr{AEDF} در داپلر گزارش شد، به علاوهٔ فلو غیرطبیعی \lr{DV}: بازگشت جریان در انتهای دیاستول.
\item انتخاب روش زایمان باید بر اساس هر بیمار باشد.
\item ولی باید بدانیم که ذخیرهٔ جفتی این جنین‌ها بسیار اندک است.
\item اگر در مانیتورینگ حین زایمان شواهد اختلال در جنین مشاهده شد: \lr{C/S}.
\item اگر نگرانی جدی دربارهٔ سلامت جنین وجود داشت: \lr{Pre-labour C/S}.
\item \lr{AEDF} یا جریان بازگشتی انتهای دیاستولی.
\end{itemize}

\subsection*{عوارض طولانی‌مدت \lr{FGR}}
\lr{DM-2}، بیماری عروق کرونری، \lr{handycap}، \lr{IQ}، ۵–۱۰ واحد کمتر.
\FGnote{املای \lr{handycap} مطابق نوشتهٔ منبع است.}

\subsection*{الگوریتم پیگیری جنین \lr{SGA}}
\FGmath{\mathrm{AC}/\mathrm{EFW}<10\mathrm{th}}

داپلر \lr{(UAD)} (شریان آمبیلیکال). شاخه‌های نمودار به ترتیب زیر هستند:

\paragraph{۱. داپلر نرمال}
\lr{AFV} (حجم مایع آمنیوتیک) بررسی می‌شود.
\begin{itemize}
\item \textbf{\lr{AFV} نرمال:} \lr{UAD, AFV, AC/EFW} هر ۲ هفته.
\item \textbf{الیگوهیدرآمنیوس:} \lr{UAD} و \lr{AFV} هر هفته؛ \lr{AC, EFW} هر ۲ هفته.
\item اگر همهٔ بررسی‌ها طبیعی \lr{NL}، زایمان در ۳۷–۳۸.
\end{itemize}

\paragraph{۲. \lr{RI} (اندکس مقاومت) بالاتر از صدک ۹۵}
\FGmath{95\mathrm{th}<\mathrm{RI}}
\begin{itemize}
\item \lr{AFV} نرمال: \lr{UAD, AFV} هر هفته؛ \lr{AC, EFW} هر ۲ هفته.
\item اگر همه \lr{NL}، زایمان در ۳۷–۳۸.
\item الیگو: \lr{CTG} اگر \FGmath{26<GA} هفته.
  \FGnote{در ابتدای این شاخه \lr{AFV} خط خورده و «الیگو» جایگزین شده است.}
\item \lr{CTG} نرمال: \lr{AFD, UAD, CTG} هفته‌ای ۲ بار؛ \lr{AC, EFW} هر ۲ هفته.
  \FGnote{در این شاخه و شاخهٔ \lr{AEDF}، اختصار حجم مایع شبیه \lr{AFD} نوشته شده؛ در شاخه‌های بالاتر \lr{AFV} است.}
\item \lr{CTG} غیرطبیعی: مشاورهٔ نوزادان، سپس زایمان.
\end{itemize}

\paragraph{۳. داپلر غیرطبیعی با \lr{AEDF}}
\begin{itemize}
\item \lr{AC/EFW} هر ۲ هفته (یادداشت حاشیه).
\item \lr{AEDF} زیر ۳۴ هفته: بدون اندیکاسیون دیگر تا زایمان؛ \lr{CTG} روزانه؛ \lr{AFD, UAD} دو بار در هفته؛ کورتون.
  \FGunclear{عبارت پس از «بدون اندیکاسیون دیگر»: «تا زایمان»؛ احتمال خوانش «تأخیر زایمان».}
\item بالای ۳۴ هفته: زایمان.
\end{itemize}

\paragraph{۴. \lr{REDF}}
بستری / \lr{CTG} روزانه / داپلر پیشرفته / استروئید / زایمان \lr{C/S}.

\clearpage
\section*{پیوست: تصاویر کامل منبع}
\begin{center}\textbf{\lr{pic01} — صفحهٔ ۱}\par\FGimage{pic01}\end{center}
\clearpage
\begin{center}\textbf{\lr{pic02} — صفحهٔ ۲}\par\FGimage{pic02}\end{center}
\clearpage
\begin{center}\textbf{\lr{pic03} — صفحهٔ ۳}\par\FGimage{pic03}\end{center}
\clearpage
\begin{center}\textbf{\lr{pic04} — صفحهٔ ۴}\par\FGimage{pic04}\end{center}
\clearpage
\begin{center}\textbf{\lr{pic05} — صفحهٔ ۵}\par\FGimage{pic05}\end{center}
\clearpage
\begin{center}\textbf{\lr{pic06} — صفحهٔ ۶}\par\FGimage{pic06}\end{center}
\FGfinish
